\section{Probabilistic Forecasting --- Quantile Regression with Split Conformal Calibration}
\label{sec:quantile}

The deterministic chapters above produce point forecasts. For a hedge-fund deployment, that is not enough: an actionable forecast also needs a calibrated estimate of its own uncertainty so that position sizing can reflect predictive confidence. This chapter builds prediction \emph{intervals} on the four representative series using LightGBM quantile regression with split conformal post-calibration, evaluated under a coverage-first rolling backtest.

\subsection{Coverage-first rolling backtest}
\label{sec:quantile:protocol}

For each of the four series the panel is split into rolling chronological folds. Each fold consists of three contiguous segments: a training segment, a 20-step calibration segment, and a 20-step validation segment. The minimum training length is 70 rows and the window expands by 20 rows per fold (one validation segment), yielding 2--6 folds per series. The protocol preserves chronology --- training points always precede calibration points, which precede validation points --- and uses canonical split semantics throughout.

For each fold three candidate methods are evaluated:

\begin{itemize}
  \item \texttt{lgb\_pair}: a standalone LightGBM at lower-quantile $\tau_{\text{lo}}$ and upper-quantile $\tau_{\text{hi}}$, scored uncalibrated.
  \item \texttt{lgb\_conformal}: the same LightGBM, post-calibrated by split conformal --- the predicted interval on the calibration slice is inflated by an additive constant so that calibration weighted PICP equals the finite-sample-corrected target $\lceil(n_{\text{cal}} + 1)\cdot \text{nominal}\rceil / n_{\text{cal}}$.
  \item \texttt{naive\_lag1}: a structural baseline using the lag-1 forecast plus a width derived from the quantiles of the in-sample residuals.
\end{itemize}

For each method we sweep all $(\tau_{\text{lo}}, \tau_{\text{hi}})$ pairs from the grid $\{0.02, 0.05, 0.10, 0.25, 0.50, 0.75, 0.90, 0.95, 0.98\}$ (lower $<$ upper) and pick the pair whose weighted calibration PICP is closest to the finite-sample target. Width acts as a tiebreaker --- between two pairs with equal coverage distance, the narrower interval wins. Across folds we then pick the dominant method per series as the one with the smallest mean coverage distance.

\subsection{Per-series leaderboard}
\label{sec:quantile:leaderboard}

\Cref{tab:quantile_leaderboard} reports the per-series summary on the stitched validation predictions, with 500-resample bootstrap intervals on weighted PICP and on the median-point skill.

\begin{table}[h]
\centering
\caption{Per-series probabilistic forecasting summary on the stitched validation predictions. ``wPICP'' is the weighted empirical coverage of the predicted 80\% interval. ``Skill (median)'' is the weighted skill score of the LightGBM median forecast. Bootstrap CIs are 500 row-resamples each.}
\label{tab:quantile_leaderboard}
\begin{tabular}{l l c c c}
\toprule
Series & Method & $n$ val & wPICP (95\% CI) & Skill of median (95\% CI) \\
\midrule
Highest total weight & lgb\_conformal & 40  & $0.60$ \,$[0.45, 0.74]$ & $\mathbf{0.61}$ \,$[0.48, 0.71]$ \\
Longest history      & lgb\_conformal & 100 & $0.82$ \,$[0.74, 0.89]$ & $0.09$ \,$[0.00, 0.35]$ \\
Most stable          & lgb\_conformal & 40  & $0.82$ \,$[0.70, 0.93]$ & $0.04$ \,$[0.00, 0.32]$ \\
Most volatile        & lgb\_conformal & 40  & $0.72$ \,$[0.59, 0.86]$ & $\mathbf{0.83}$ \,$[0.76, 0.89]$ \\
\bottomrule
\end{tabular}
\end{table}

The median LightGBM produces substantively positive skill on the most-volatile and highest-weight series ($0.83$ and $0.61$ respectively), and near-zero skill on the longest-history and most-stable series --- consistent with the regimes documented in \cref{sec:eda:repseries}. The 80\% intervals achieve weighted PICP in the range $0.60$--$0.82$. Three of the four series have bootstrap CIs that include the nominal $0.80$ target; only the highest-weight series under-covers significantly (CI $[0.45, 0.74]$, excluding $0.80$). The longest-history and most-stable series land essentially on target.

\subsection{Method comparison}
\label{sec:quantile:methods}

Although \texttt{lgb\_conformal} is the dominant method for every series under our calibration-distance selection rule, the validation behaviour of all three methods is informative. \Cref{fig:quantile_method_comparison} plots weighted PICP and median-point skill for each (series, method) pair.

\begin{figure}[h]
\centering
\includegraphics[width=\linewidth]{05_method_comparison}
\caption{Method comparison across the four representative series. Left: weighted PICP at the nominal 80\% target (red dashed line). Right: skill of the LightGBM median forecast (the median is identical across all three methods, so the right panel reads as the underlying point-forecast skill on the validation segment).}
\label{fig:quantile_method_comparison}
\end{figure}

Two observations:

\begin{itemize}
  \item Conformal calibration consistently lifts the validation PICP toward the nominal target compared to uncalibrated \texttt{lgb\_pair}: on the longest-history series PICP rises from $0.78$ (pair) to $0.82$ (conformal), on the most-stable series from $0.74$ to $0.82$.
  \item On the most-volatile series, \texttt{naive\_lag1} actually achieves the best validation PICP ($0.83$) among the three, with a much narrower width than conformal LightGBM ($185$ vs $427$ in the original target units). This is consistent with the legacy project artefacts where \texttt{naive\_lag1} won on the volatile and high-weight series for the same reason: when the lag-1 persistence is strong, a structural baseline with quantile-derived widths is hard to beat.
\end{itemize}

\subsection{Calibration plot}
\label{sec:quantile:calibration}

\Cref{fig:quantile_calibration} plots empirical wPICP against the nominal coverage target for each series's dominant method. Three of four series land within a few percentage points of the diagonal --- visually confirming that conformal calibration is delivering close to the promised marginal coverage.

\begin{figure}[h]
\centering
\includegraphics[width=0.7\linewidth]{05_calibration_plot}
\caption{Empirical weighted PICP versus the nominal $80\%$ target, one point per series for the dominant method. Diagonal dashed line is perfect calibration. The longest-history and most-stable series sit very close to the diagonal; the most-volatile and highest-weight series under-cover slightly.}
\label{fig:quantile_calibration}
\end{figure}

\subsection{Coverage by row weight}
\label{sec:quantile:coverage_weight}

Because the official metric is sample-weighted, coverage on high-weight rows matters more than coverage on low-weight rows. \Cref{fig:quantile_coverage_weight} bins the validation rows by weight decile, per series, and shows empirical PICP per bin.

\begin{figure}[h]
\centering
\includegraphics[width=\linewidth]{05_coverage_by_weight}
\caption{Empirical PICP by weight decile, one panel per series. Red dashed line is the nominal $80\%$ target. The validation segments are small ($n = 40$--$100$) so per-decile estimates are noisy, but the trend across deciles indicates whether high-weight rows receive comparable coverage to low-weight rows.}
\label{fig:quantile_coverage_weight}
\end{figure}

The coverage tends to fluctuate around the nominal target without a systematic decline at the high-weight extreme, which is the property a deployable interval needs --- the metric is dominated by high-weight rows and they are not systematically under-covered.

\subsection{Fan charts on the four representative series}
\label{sec:quantile:fan_charts}

For each of the four series we render a full-width fan chart with the historical training segment, the validation truth, the LightGBM median forecast, and the calibrated $80\%$ interval. Each chart carries an annotation box reporting weighted PICP, raw PICP, weighted MPIW, raw MPIW, the median's point skill, the number of folds used, the dominant method, and the selected quantile pair. The interval band is the headline output of this chapter --- a deployable forecast that is not just a point estimate but a calibrated range with quantified uncertainty.

\begin{figure}[h]
\centering
\includegraphics[width=\linewidth]{05_fan_longest_history}
\caption{Longest-history series. Coverage is well-calibrated (wPICP $0.82$); point skill is near zero, which is honest --- the data has no exploitable feature signal at this resolution.}
\label{fig:fan_longest}
\end{figure}

\begin{figure}[h]
\centering
\includegraphics[width=\linewidth]{05_fan_highest_total_weight}
\caption{Highest-total-weight series. The $80\%$ interval covers most observations; point skill of the median is $0.61$, the second-strongest empirical result in the project. Coverage slightly under-shoots the nominal $0.80$ at $0.60$.}
\label{fig:fan_highest}
\end{figure}

\begin{figure}[h]
\centering
\includegraphics[width=\linewidth]{05_fan_most_volatile}
\caption{Most-volatile series. The headline visual of the chapter: median forecast tracks the validation trajectory closely (skill $0.83$), the calibrated band hugs the truth, and weighted PICP is $0.72$ on a 40-row validation segment.}
\label{fig:fan_volatile}
\end{figure}

\begin{figure}[h]
\centering
\includegraphics[width=\linewidth]{05_fan_most_stable}
\caption{Most-stable series. The series itself is near-flat at scale $10^{-4}$, so the median forecast carries no signal; the calibrated band still covers most observations (wPICP $0.82$), which is the correct behaviour: the interval acknowledges its own uncertainty rather than committing to a confident wrong point estimate.}
\label{fig:fan_stable}
\end{figure}

The most-volatile panel is the visual headline: the median tracks the validation trajectory closely (skill $0.83$) and the band covers the true values most of the time. The highest-weight panel operates at scale $10^{-4}$ but the band still hugs the truth and captures most of the swings. The longest-history and most-stable panels show the band capturing most observations even when the median forecast itself adds little signal --- exactly the right behaviour: an honest interval acknowledges that point prediction is hard there and reports its uncertainty appropriately.

\subsection{What this chapter establishes}
\label{sec:quantile:conclusions}

Three points to carry forward:

\begin{enumerate}
  \item \textbf{Calibrated intervals are deployable.} On three of four series the empirical weighted PICP is within bootstrap-CI of the nominal $80\%$ target. The under-coverage on the highest-weight series is the one place where exchangeability between calibration and validation appears to break.
  \item \textbf{The structural baseline is competitive on volatile series.} Naive lag-1 with quantile-derived widths achieves narrower, better-calibrated intervals on the most-volatile series than conformal LightGBM. This is a real result, not a methodological problem with conformal: the cumulation property already makes lag-1 a strong forecast at the volatile-series scale, so a structural method that exploits it directly beats a regression model that has to learn the same property.
  \item \textbf{Probabilistic forecasting closes the loop into actions.} Section~\ref{sec:actions} discusses how a calibrated interval --- not just a point estimate --- is what's needed to translate the model output into deployable position sizes. The chapter above is the missing link between the deterministic models in the previous chapters and the actionable downstream use.
\end{enumerate}
